% ============================================================================
% CAPÍTULO 16 — Testes automatizados
% Status: texto completo (rodada editorial 2)
% ============================================================================
\chapter{Testes Automatizados}
\label{cap:testes}

\begin{caixaconceito}
Testar não é opcional: é a rede de segurança que permite refatorar sem medo,
documenta o comportamento e separa código ``que funciona'' de código
\emph{confiável}. A stack de mercado: \textbf{Vitest} para unit/componente e
\textbf{Playwright} para E2E \cite{vitest-docs,playwright-docs}.
\end{caixaconceito}

Existe um momento em todo projeto em que o código ``funciona'' mas ninguém
ousa mexer nele. A funcionalidade nova quebra a antiga; o bug corrigido
ressuscita três semanas depois; o deploy vira um evento de risco. Esse momento
não é sobre escrever código — é sobre \emph{confiança}. Testes automatizados
são a ferramenta que transforma ``espero que funcione'' em ``sei que
funciona''.

Este capítulo não trata de ``técnica de teste'' como abstração: trata da
suíte que você vai escrever de verdade, no projeto real do SkillHub, com as
ferramentas que o mercado TypeScript adotou em massa — Vitest\index{Vitest} para testes
rápidos e Playwright\index{Playwright} para fluxos completos no navegador.

Ao final deste capítulo, você será capaz de:
\begin{objetivos}
  \item Aplicar a pirâmide de testes\index{pirâmide de testes} (unitário $\rightarrow$ integração $\rightarrow$ E2E);
  \item Escrever testes unitários\index{testes unitários} com Vitest (funções e componentes React);
  \item Testar integrações com banco real ou em memória;
  \item Escrever testes E2E com Playwright (fluxos reais do usuário);
  \item Medir cobertura\index{cobertura} e escrever testes que valem a pena;
  \item Entregar o projeto do capítulo: a suíte completa do SkillHub.
\end{objetivos}

\section{A pirâmide de testes}

A pirâmide organiza os testes por \emph{custo} e \emph{quantidade}:

\begin{itemize}[leftmargin=*, itemsep=0.4em]
  \item \textbf{Unitários} (base, muitos): funções e componentes isolados — rápidos;
  \item \textbf{Integração} (meio): módulo + banco + HTTP — verificam o encaixe;
  \item \textbf{E2E} (topo, poucos): fluxo completo no navegador — caros, mas validam o produto inteiro.
\end{itemize}

\begin{caixaconceito}
A pirâmide existe porque o custo (tempo, manutenção, flakiness) cresce de
baixo para cima. Priorize muitos testes unitários baratos e poucos E2E
críticos. Testes E2E que testam tudo viram pesadelo de manutenção.
\end{caixaconceito}

\begin{caixadica}
Uma forma útil de pensar: um teste unitário responde ``esta função faz o que
promete?'', um teste de integração responde ``essas peças se encaixam?'', e um
E2E responde ``o usuário consegue concluir a tarefa?''. As três perguntas são
diferentes — e todas precisam de resposta.
\end{caixadica}

\subsection{TDD: teste primeiro, ou depois?}

Test-driven development (escrever o teste antes da implementação) é uma
ferramenta poderosa de \emph{design}: ele força você a pensar na interface
antes da implementação. Mas não é dogma. No dia a dia, o padrão mais comum é
híbrido: testes críticos (regras de negócio, validação, bugs corrigidos) são
escritos com TDD\index{TDD}; testes de cobertura de utilitários simples vêm depois. O
importante não é a ordem — é \emph{existirem} e rodarem a cada mudança.

\section{Vitest: testes unitários rápidos}

O Vitest é o sucessor moderno do Jest para o ecossistema Vite/TypeScript:
mesma API familiar (\texttt{describe}/\texttt{it}/\texttt{expect}), mas com
transformação nativa de TypeScript e modo watch instantâneo:

\begin{lstlisting}[style=livro, language=TypeScript,
  caption={Teste unitário com Vitest.}, label={list:teste-vitest}]
import { describe, it, expect } from "vitest";
import { formatarMoeda, validarEmail } from "./utils";

describe("formatarMoeda", () => {
  it("formata valores em reais", () => {
    expect(formatarMoeda(1234.5)).toBe("R$ 1.234,50");
  });
});

describe("validarEmail", () => {
  it("aceita e-mails válidos", () => {
    expect(validarEmail("ana@exemplo.com")).toBe(true);
  });

  it("rejeita e-mails sem arroba", () => {
    expect(validarEmail("anaexemplo.com")).toBe(false);
  });
});
\end{lstlisting}

\begin{lstlisting}[style=livro, language=Bash,
  caption={Executando a suíte.}, label={list:teste-run}]
npx vitest run        # uma vez
npx vitest            # modo watch
npx vitest run --coverage   # relatório de cobertura
\end{lstlisting}

\begin{caixadica}
Escolha \textbf{assertions} pelo que elas verificam: \texttt{toBe} compara
identidade (primitivos), \texttt{toEqual} compara estrutura (objetos e
arrays), \texttt{toContain} verifica presença em coleções, e
\texttt{toThrow} valida exceções. Usar o matcher errado é a causa nº 1 de
testes que passam sem testar nada.
\end{caixadica}

\subsection{Mocks: isolando o que não é seu}

Quando a função testada depende de algo externo (banco, HTTP, relógio), você
tem duas escolhas: subir a dependência de verdade (teste de integração) ou
substituí-la por um \emph{mock} (teste unitário). O Vitest traz
\texttt{vi.fn()}, \texttt{vi.mock()} e \texttt{vi.useFakeTimers()}:

\begin{lstlisting}[style=livro, language=TypeScript,
  caption={Mock de fetch e de relógio.}, label={list:teste-mock}]
import { vi, it, expect } from "vitest";
import { buscarUsuario } from "./api";

it("trata falha de rede com fallback", async () => {
  vi.spyOn(globalThis, "fetch").mockResolvedValueOnce(
    new Response("{}", { status: 500 }),
  );

  await expect(buscarUsuario("ana")).resolves.toBeNull();
});

it("usa o tempo atual no log", () => {
  vi.useFakeTimers();
  vi.setSystemTime(new Date("2026-08-08T12:00:00Z"));
  // ... código que depende de Date.now()
  vi.useRealTimers();
});
\end{lstlisting}

\begin{caixaarmadilha}
Mocks\index{mocks} mentem. Um mock de \texttt{prisma.servico.create} não valida schema,
não aplica constraints e não aciona triggers — ele devolve exatamente o que
você programou. Use mocks para \emph{isolar}, não para \emph{simular o banco}:
regras de banco merecem teste de integração real.
\end{caixaarmadilha}

\section{Testando componentes React}

Use a \textbf{Testing Library}: teste o que o usuário \emph{vê e faz}, não os
detalhes internos:

\begin{lstlisting}[style=livro, language=TypeScript,
  caption={Teste de componente com Testing Library.}, label={list:teste-react}]
import { render, screen, fireEvent } from "@testing-library/react";
import { Contador } from "./Contador";

it("incrementa ao clicar", () => {
  render(<Contador />);
  const botao = screen.getByRole("button", { name: /cliques/i });
  fireEvent.click(botao);
  fireEvent.click(botao);
  expect(screen.getByText("Cliques: 2")).toBeInTheDocument();
});
\end{lstlisting}

\begin{itemize}[leftmargin=*, itemsep=0.4em]
  \item \textbf{getByRole} — encontra pelo papel acessível (o mesmo que leitores de tela usam);
  \item \textbf{getByLabel} — encontra campos pelo \texttt{<label>} (a mesma forma que o usuário identifica);
  \item \textbf{getByText} — conteúdo visível;
  \item \textbf{fireEvent} vs \textbf{userEvent} — \texttt{userEvent} simula interações reais (digitação, foco) e é mais fiel; prefira-o.
\end{itemize}

\begin{caixaconceito}
Se um teste depende de \texttt{className} ou do HTML interno de um componente,
ele quebra quando o design muda — e você perde a confiança na suíte. Testar
por papel acessível é testar o contrato com o usuário, que muda muito menos
que o markup.
\end{caixaconceito}

\section{Testes de integração com banco}

O teste de integração sobe a dependência real — banco, HTTP — e verifica o
encaixe. Para o Prisma, o padrão é um banco PostgreSQL de teste isolado:

\begin{lstlisting}[style=livro, language=TypeScript,
  caption={Integração: repositório contra banco real (Docker).},
  label={list:teste-integracao}]
import { describe, it, expect, beforeAll, afterAll } from "vitest";
import { prisma } from "./db";

describe("repositorio de servicos", () => {
  beforeAll(async () => {
    await prisma.$executeRaw`TRUNCATE servicos RESTART IDENTITY CASCADE`;
  });

  afterAll(async () => {
    await prisma.$disconnect();
  });

  it("cria e busca um serviço", async () => {
    const criado = await prisma.servico.create({
      data: { titulo: "Aula de React", precoCentavos: 5000 },
    });
    const encontrado = await prisma.servico.findUnique({ where: { id: criado.id } });
    expect(encontrado?.titulo).toBe("Aula de React");
  });
});
\end{lstlisting}

\begin{caixadica}
Para integração com banco, rode o PostgreSQL em Docker (capítulo~\ref{cap:docker})
e use um banco de teste separado (\texttt{DATABASE\_URL} apontando para
\texttt{...\_test}). Isolamento total entre testes = resultados confiáveis.
\end{caixadica}

\begin{caixadica}
Em vez de \texttt{TRUNCATE} manual, considere rodar \texttt{prisma migrate
reset} ou \texttt{prisma db push} no banco de teste no \texttt{beforeAll} da
suíte. Isso garante que o schema de teste esteja sempre sincronizado com as
migrações — e pega erros de migração antes do deploy.
\end{caixadica}

\section{Playwright: testes E2E no navegador real}

O Playwright abre um navegador real (Chromium, Firefox, WebKit), executa o
fluxo como um usuário e verifica o resultado na tela:

\begin{lstlisting}[style=livro, language=TypeScript,
  caption={E2E: o fluxo completo de um pedido.}, label={list:teste-e2e}]
import { test, expect } from "@playwright/test";

test("profissional anuncia um serviço", async ({ page }) => {
  await page.goto("/");
  await page.getByRole("button", { name: "Entrar" }).click();
  await page.getByLabel("E-mail").fill("ana@exemplo.com");
  await page.getByLabel("Senha").fill("senha-forte-123");
  await page.getByRole("button", { name: "Entrar" }).click();

  await page.goto("/painel");
  await page.getByRole("button", { name: "Novo serviço" }).click();
  await page.getByLabel("Título").fill("Aula de Next.js");
  await page.getByRole("button", { name: "Publicar" }).click();

  await expect(page.getByText("Aula de Next.js")).toBeVisible();
});
\end{lstlisting}

\begin{caixaconceito}
O Playwright espera por \emph{condições}, não por \emph{tempo}: locators
(\texttt{getByRole}, \texttt{getByLabel}) + \texttt{expect} com auto-wait
eliminam os \texttt{sleep} que tornam testes frágeis. Teste como usuário: por
papel, rótulo e texto visível.
\end{caixaconceito}

\begin{caixadica}
Crie dados de teste com \textbf{API} (não pela UI): no \texttt{beforeAll},
cadastre usuários e serviços via Route Handlers ou Prisma direto. A UI fica
livre para testar o que ela deve testar — o fluxo — em vez de perder tempo
repetindo o cadastro em todo teste.
\end{caixadica}

\section{Cobertura: métrica, não meta}

Cobertura mede \emph{quantas linhas rodaram}, não \emph{se o comportamento está
correto}. Use-a para encontrar código morto e lacunas, não como número a
inflar. Uma linha de negócio crítica testada vale mais que 100\% de cobertura
de funções triviais.

\begin{caixaconceito}
\textbf{O teste que vale a pena} é aquele que falha quando o comportamento
muda. Se você pode remover um teste e nada quebra, ele não estava protegendo
nada — ele estava só decorando o relatório de cobertura.
\end{caixaconceito}

% ============================================================================
\section{Projeto do capítulo: a suíte do SkillHub}
\label{sec:projeto16}

\begin{caixaprojeto}
\textbf{Objetivo:} criar a suíte de testes completa do marketplace SkillHub
(capítulo~\ref{cap:fullstack}).

\textbf{Critérios de aceite:}
\begin{itemize}[leftmargin=*, itemsep=0.2em]
  \item \textbf{Unitários}: utilitários (formatação, validação, base62) e componentes isolados;
  \item \textbf{Integração}: repositório Prisma contra banco de teste + rotas HTTP;
  \item \textbf{E2E}: 3 fluxos — cadastro/login, anunciar serviço, contratar serviço;
  \item Testes de regressão do RBAC (não-dono recebe \texttt{403});
  \item Scripts no \texttt{package.json}: \texttt{test}, \texttt{test:e2e}, \texttt{coverage};
  \item CI executando a suíte (capítulo~\ref{cap:cicd});
  \item Cobertura $\geq$ 70\% nas funções de negócio, documentada.
\end{itemize}
\end{caixaprojeto}

\subsection{Passo a passo}

\begin{passos}
  \item \textbf{Instalação}: adicione Vitest + Testing Library ao projeto e crie
  o script \texttt{test} no \texttt{package.json};
  \item \textbf{Unitários}: teste \texttt{formatarMoeda}, \texttt{validarEmail}
  e o codificador/decodificador (use property-based com \texttt{fast-check}
  como desafio extra);
  \item \textbf{Componentes}: teste o \texttt{FormularioDeServico} — erro de
  validação visível, estado \texttt{pendente} e sucesso;
  \item \textbf{Integração}: suba o PostgreSQL de teste em Docker, rode as
  migrações e teste o repositório Prisma + as Server Actions (sucesso e
  falhas de RBAC);
  \item \textbf{E2E}: configure o Playwright, crie os 3 fluxos e rode contra o
  servidor de desenvolvimento com banco de teste;
  \item \textbf{CI}: faça o GitHub Actions rodar unit + integração + E2E em
  cada pull request (capítulo~\ref{cap:cicd});
  \item \textbf{Documentação}: registre como rodar cada camada e o relatório
  de cobertura no \texttt{README} do projeto.
\end{passos}

\section{Exercícios de fixação}

\begin{exercicios}
  \item Desenhe a pirâmide de testes e explique por que ela existe.
  \item Escreva testes unitários para \texttt{validarEmail} (4 casos) e \texttt{formatarMoeda}.
  \item Qual a diferença entre \texttt{toBe} e \texttt{toEqual} no Vitest?
  \item Escreva um teste de componente: um botão que mostra ``Carregando...'' enquanto pendente.
  \item Por que locators por \texttt{role} são melhores que por \texttt{class}?
\end{exercicios}

\section{Desafios}

\begin{desafios}
  \item \textbf{Teste de Server Action}: teste \texttt{criarServico} com validação Zod (sucesso, título curto, preço negativo).
  \item \textbf{Mock de fetch}: teste um componente que busca dados com \texttt{vi.mock} e estados de erro.
  \item \textbf{Snapshot visual}: use o screenshot do Playwright para comparar o design system antes/depois de uma mudança de tema.
\end{desafios}

\section{Exercícios de nível sênior}

\begin{seniores}
  \item \textbf{Testes determinísticos}: elimine flakiness de um teste E2E que falha 1 em 10 vezes (dica: isolamento de dados, locators, retries) e documente a causa raiz.
  \item \textbf{Teste de carga}: use K6 para 100 usuários simultâneos na rota de busca e defina um orçamento de performance (p95 $<$ 300ms).
  \item \textbf{Property-based testing}: use \texttt{fast-check} para gerar 1000 URLs aleatórias e validar o par codificar/decodificar do encurtador (capítulo~\ref{cap:node}).
  \item \textbf{Mutação}: rode o \texttt{stryker} e elimine mutantes sobreviventes nas funções de negócio críticas.
\end{seniores}

\section{Armadilhas comuns}

\begin{itemize}[leftmargin=*, itemsep=0.4em]
  \item Testar implementação em vez de comportamento (testes que quebram a cada refactor);
  \item \texttt{sleep} em E2E (testes frágeis — use auto-wait);
  \item Cobertura como meta em vez de ferramenta;
  \item Esquecer de testar os caminhos de erro (401, 403, 404, 500);
  \item Banco de teste compartilhado com dados vazando entre testes;
  \item Testes que dependem de ordem de execução ou de estado global (torne cada teste independente);
  \item E2E testando detalhes internos da UI (ex.: \texttt{class}) em vez do comportamento visível.
\end{itemize}

\section{Resumo do capítulo}

\begin{itemize}[leftmargin=*, itemsep=0.3em]
  \item Pirâmide: muitos unitários, alguns de integração, poucos E2E;
  \item Vitest para unit/componente; Testing Library para comportamento;
  \item Integração com banco real isolado (Docker);
  \item Playwright para fluxos reais com locators acessíveis;
  \item Cobertura é ferramenta, não meta;
  \item Projeto entregue: suíte completa do SkillHub.
\end{itemize}

\section{Referências oficiais para aprofundar}

\begin{itemize}[leftmargin=*, itemsep=0.3em]
  \item Documentação do Vitest: \url{https://vitest.dev}
  \item Testing Library: \url{https://testing-library.com}
  \item Documentação do Playwright: \url{https://playwright.dev}
  \item Jest (alternativa madura): \url{https://jestjs.io}
\end{itemize}
